\section{Discussion and Conclusion}
\label{sec:discussion}

\macaron{} is our first release of an agent model family designed around
continual learning and collective intelligence rather than around a single
monolithic checkpoint~\citep{macaron_v1_blog2026}. The architectural spine is
\molshort{}: a frozen base plus a small set of specialist LoRA
adapters~\citep{lora2022,mol_harness2026}, orchestrated by a Proxy-mediated route hop
whose label is emitted by L0, and served through a harness we treat as a first-class training target
~\citep{ui4a2026,wu2026replharnesses}. The training spine is \mindforge's
recursive self-improvement loop~\citep{macaron_v1_preview2026}, running against
the same harness that ships in production, over an infrastructure stack
(\mint~\citep{lu2026announcing}, \longstraw{}~\citep{zhou2026longstrawlongcontextrl2m})
built to support the multi-million-token operating points reported in
Section~\ref{sec:infra_longstraw}.
The evaluation spine is Personal Intelligence, expressed through \chatbench{}
and \livingbench.

\subsection{Adaptation and Collaboration, Revisited}
\label{sec:disc_bets}

Section~\ref{sec:introduction} framed \macaron{} around two bets, and it is worth restating them in light of what the rest of the report actually delivers.

\emph{Adaptation} is not the property of a single model; it is the property of the loop the model sits in~\citep{silver_sutton_experience_2025,yao_second_half_2025}. The RSI cycle (Section~\ref{sec:rsi}) generates harder tasks, audits trajectories under a versioned harness configuration, evaluates candidate HCPs, and only then uses selected trajectories to update weights. The design description concerns this loop; the released checkpoint is one snapshot and does not demonstrate improvement across generations.

\emph{Collaboration} is not a multi-agent framework layered on top of a model; it is an architectural affordance. Because the base is shared and adapters are portable, the registry can admit specialists trained by different teams or personalized for different users on the same runtime~\citep{macaron_v1_blog2026,lu2026announcing}. The Proxy's routing interface is the proposed interoperability contract; this release tests only the four shipped specialists.

Neither bet is settled by \macaron{}. The design is inspectable in the released harness (\url{https://github.com/MindLab-Research/Mixture-of-LoRA-Harness}) and \venti{} weights (\url{https://huggingface.co/mindlab-research/Macaron-V1-Venti}); cross-generation and independently trained specialist evaluations remain work for the next release.

\subsection{Limitations and Focus Directions}
\label{sec:disc_limitations}

We frame the limitations of this release around five directions that most directly shape the next generation. Each is a place where the current evidence falls short of the system goal it serves.

\paragraph{Scaling recursive self-iteration.}
What we release is one snapshot of an RSI loop, not a system that has demonstrably compounded over many generations. The released \venti{} snapshot cannot, by itself, distinguish a compounding RSI effect from a single round of self-generated-data training, and cross-generation lift is not yet measured. Two scaling questions follow. First, how to keep the model learning and iterating on top of the competence it has already reached, rather than re-deriving behavior each generation; self-generated tasks tend to converge on the shapes the current policy finds discoverable, which is not the shape of real user behavior, and the loop can optimize itself into a local mode~\citep{macaron_v1_blog2026}. Second, how to run RSI at larger scale: more tasks, longer contexts, and the multi-million-token rollouts the \longstraw{} infrastructure makes executable but has not yet been driven through a full learning curve (Section~\ref{sec:infra_longstraw}). Both are prerequisites for the continual-learning bet to compound rather than reset.

\paragraph{Scenario richness.}
The Personal Intelligence benchmarks (\chatbench, \livingbench) and the general-capability suite cover a real but narrow slice of the situations a personal agent meets. \chatbench{} test items come from de-identified product conversations, and \livingbench{} seeds from real failure taxonomies, but simulator quality is a moving target: any residual mismatch between the simulator and a real user creates an evaluation--deployment gap, and we observe qualitative character-stability degradation after multiple preference-drift events in some simulated sessions. The next generation requires more diverse scenarios and a quantified stability analysis so that the trajectory distribution optimized by the RSI loop better reflects deployment.

\paragraph{Evaluation scope and provenance.}
\chatbench{} and \livingbench{} target the same Personal Intelligence
distribution that informs the RSI loop, so they measure the behaviors this
release is designed to improve. Every unstarred model within a benchmark row is
evaluated on the same task set and benchmark-specific protocol, which supports
direct comparison on that benchmark. The internal suites remain focused in
size and scope, and several starred general-benchmark values come from public
leaderboards; we retain those provenance markers and use the imported values as
context. Broader public test releases and additional judge calibration will
extend the coverage of future evaluations.

\paragraph{Data governance, reproducibility, and safety scope.}
The internal evaluations include de-identified product conversations and traffic,
but this report does not document the consent or opt-in basis for research use,
the de-identification procedure and residual re-identification audit, retention
and access controls, or an ethics-review determination. The released artifacts
document the model--harness interfaces and selected operating points, but this
report also does not provide a complete per-specialist training specification or
a standalone safety and red-team evaluation. These are material limitations for
auditing the data pipeline, reproducing the reported checkpoints, and assessing
deployment in high-stakes personal-assistance settings. We therefore interpret
the present results as a systems characterization, not as evidence that the
release is suitable for safety-critical use.

\paragraph{Emergence of collective intelligence.}
\molshort{} is designed for collective intelligence: specialists trained by different teams or personalized for different users are composed on the same shared base (Section~\ref{sec:mol_collective}). This release has not yet demonstrated it. The current tests exercise routing, per-specialist conversation views, and adapter registration for the four shipped specialists; they do not establish robust switching under broad workloads or show that a population trained by different teams or users produces capability beyond any constituent specialist. Whether collective intelligence is an emergent property of the architecture, rather than an affordance, is the central open question for the next generation.

\subsection{Roadmap}
\label{sec:disc_roadmap}

We keep the roadmap short and concrete.

\paragraph{More specialists.}
\molshort{} is set up for additive capability. The near-term specialists we are training target areas not covered by the first four, including domain-specific research and non-English long-form. Decision-support domains require a dedicated safety evaluation before any release. Each specialist would ship as an adapter registration rather than as a base retrain, leaning on \mint's measured addressable catalog~\citep{lu2026announcing}.

\paragraph{More rendering targets for \uifora.}
Flutter and native mobile are the next renderers~\citep{ui4a2026}. The protocol does not change; adding a target ships a small renderer and a curated component import layer.

\paragraph{Third-party adapters and personalization.}
Collective intelligence is a claim that has to be shown, not just designed for. The next milestone is a public composition path for third-party adapters and a personalization path for user-specific adapters, both on the same runtime as the shipping specialists~\citep{macaron_v1_blog2026}.

\paragraph{Deeper integration with the artifact loop.}
Real interactions through Macaron Artifacts feed the RSI loop with the trajectory distribution we most want to optimize against. Closing this loop tightly, with clear consent and opt-in, is what makes the continual-learning story compound.

\subsection{Conclusion}
\label{sec:disc_conclusion}

Building a growable agent model is not the same problem as building a strong
general model. It requires the environment, harness, training loop, and model
architecture to evolve as separately versioned but jointly evaluated layers.
\macaron{} is our first implementation of that decomposition: a frozen shared
base carrying general capability, LoRA specialists carrying differentiated
behavior, a harness carrying tools and context conventions, and an RSI loop
carrying the revision process. The current results document execution checks for parts of this stack
and one model snapshot; they leave longitudinal continual-learning gains and
beneficial composition across independently trained adapter populations as open
empirical questions.
